%!TeX root=../tese.tex
%("dica" para o editor de texto: este arquivo é parte de um documento maior)
% para saber mais: https://tex.stackexchange.com/q/78101

\chapter{Fundamentação}

\section{Escalonamento em lote em HPC}

Sistemas de computação de alto desempenho são complexos por natureza, com cada nó sendo composto
por inúmeras unidades de processamento, sejam CPUs ou GPUs, além dos dispositivos de memória (RAM, SSDs)
e dispositivos de rede necessários para habilitar e suportar o enorme volume de computações simultâneas
nesses nós. Advinda dessa complexidade, os usuários finais desses sistemas enviam os seus trabalhos para
um software intermediário, o Sistema de Gerenciamento de Trabalhos e Recursos (do inglês Resource and Job
Management System, ou RJMS), que pode ser resumido como o escalanador do sistema, observado que sua função
principal é decidir em que máquinas e em que momento um dado trabalho computacional deve ser executado. 
Esses trabalhos, ou \emph{jobs}, nada mais são do que aplicações computacionais, como programas ou scripts,
em conjunto com a quantidade de recursos requisitados, tempo de submissão, e outras informações pertinentes
para a execução do mesmo.

Para realizar tal função, o escalanador deve então decidir o modo para como compartilhar seus recursos disponíveis
entre os jobs requisitantes em um certo momento. Os dois esquemas principais de compartilhamento são o \emph{time-sharing},
onde um mesmo recurso é compartilhado entre diversos jobs, e com esses tendo a posse dos recursos por 
pequenas fatias de tempo; e o \emph{space-sharing}, onde um recurso é atribuído exclusivamente a um job até a finalização do mesmo. 
Por exemplo, com time-sharing, uma mesma unidade de processamento teria a posse das suas 
threads distribuída entre vários jobs por um curto período de tempo, enquanto com space-sharing
todas as threads estariam no controle do mesmo job por todo seu período de execução. Jobs que possuem
interação com o usuário durante sua execução (interativos) são preferívelmente executados em máquinas ou
processadores em um esquema de time-sharing, enquanto para jobs que rodam sem a necessidade de interação com 
o usuário (batch) se opta por um esquema de space-sharing.

Jobs são caracterizados em relação a duas dimensões espaciais: o seu comprimento, que diz respeito
ao tempo de execução decorrido de um job; e sua largura, que caracteriza um job a partir da quantidade
de recursos requisitados por tal. Sobre a largura, é possível classificar um job em uma de quatro
categorias: \textbf{rígidos}, onde o número de recursos que o escalonador atribui a um job é exatamente
igual ao número de recursos requisitados pelo usuário, durante toda a execução do job; \textbf{moldável},
onde o escalonador por si mesmo atribui um número de recursos a um certo job, fixo durante toda a execução
do mesmo; \textbf{evolutivo} onde o usuário desenvolve o job de uma maneira que a quantidade de recursos
necessários muda durante a sua execução, cabendo ao escalonador satisfazer tais atribuições; ou 
\textbf{maleável}, onde o escalonador por si próprio atribui um número de recursos e também controla
possíveis incorporações ou liberações de recursos.

A capacidade de estimar o comprimento de um job é de grande importância para a maioria das políticas 
de escalonamento. Apesar de alguns trabalhos na área apresentarem soluções quanto a estimação automática
do tempo de execução dos jobs por parte do próprio escalonador, na prática a maioria dos sistemas acaba
demandando que os usuários submitam o tempo de execução que eles esperam que o job tome, nesse caso 
funcionando como um tempo de execução máximo, ou \textit{walltime}. Apesar dos usuários possuírem a tendência
de superestimar esse tempo de execução máximo, visto que a subestimação do walltime levaria a finalização
forçada do job, é a solução utilizada na prática, inclusive com relatos alegando que essa superestimação
acaba sendo benéfica de acordo com a política de escalonamento utilizada. Esse tempo máximo declarado faz parte
do modelo de job "rígido" dominante em sistemas de HPC em produção e utilizado por esse trabalho, 
onde os jobs são executados por máquinas em esquema de space-sharing e com largura rígida.  

A grande maioria dos escalonadores em produção seguem uma lógica de enfileiramento, atribuindo recursos
disponíveis em determinado momento para jobs que estão esperando em uma fila, de acordo com a própria 
ordenação dessa fila. Diversos critérios de ordenação da fila são utilizados na prática, como pelo tempo
de chegada na fila, pelo tempo de execução máxima do job, seja crescente ou decrescente, pelo número de
recursos requisitados. O critério do tempo de chegada do job na fila (ou \textit{First-Come-First-Serve}, 
FCFS) é um dos mais utilizados na prática, dada a simplicidade da sua implementação e explicabilidade. 
Entretanto, esse critério aplicado sozinha corre o risco de causar a chamada fragmentação na máquina, 
com uma parte considerável dos nós do sistema tornando-se ociosos enquanto o job na posição mais 
prioritária da fila, a cabeça, espera pela liberação dos recursos necessários para sua execução.

Desenvolvidor para integrar o escalonador EASY (de Extensible Argonne Scheduling sYstem), gerenciador 
do IBM SP no Laborátório Nacional de Argonne, o \emph{backfilling} é um mecanismo que visa lidar com os
nós ociosos do sistema devidos a bloqueio pela cabeça da fila, sendo altamente implementado nos 
sistemas em produção. A partir do momento em que o job na cabeça da fila não consegue ser executado
dada a falta de recursos em relação aos seus requisitos, o escalonador começa a varrer os jobs posteriores
na fila. Apoiado nos tempos de execução máximo estimados pelos usuários, o escalonador checa se um 
job posterior na fila pode ou não ser executado no momento atual. Para isso, ele deve requisitar uma
quantidade de recursos menor ou igual a disponível naquele momento, e seu deslocamento não deve atrasar
a execução do job na cabeça da fila. 

O EASY backfilling garante que a o job na cabeça da fila não terá sua execução adiada
por jobs posteriores a ele caso os nós ociosos sejam ocupados, porém o mesmo não pode ser 
ditos dos outros jobs. Por isso, essa variação também é conhecida como backfilling \emph{agressivo},
em contraste com versões mais conservadoras, que garantem que um certo número de jobs posteriores
ao job prioritário ou mesmo nenhum job da fila tenha sua execução delegada por conta do mecanismo
de backfilling. No fim das contas, o EASY é o baseline consolidado, pensando tanto nos
escalonadores dos sistemas reais quanto na literatura de políticas de escalonamento.
